\documentclass{article}
\usepackage{iclr2026_conference,times}

\input{math_commands.tex}

\usepackage{amsmath,amssymb,amsthm,bm,mathtools}
\usepackage{booktabs}
\usepackage{graphicx}
\usepackage{enumitem}
\usepackage{microtype}
\usepackage{hyperref}
\usepackage{url}

% -----------------------------------------------------------------------------
% ICLR compact manuscript skeleton.
% This is a separate, author-facing alternative to main_manuscript_skeleton.tex.
% Bullets mark paragraph content; proofs and extended derivations are deferred.
%
% Main-text page budget (references and appendix excluded):
% Abstract + introduction                         1.25 pages
% Setup + evaluation table                        0.75 pages
% Pixel ridge DE + CV                             1.50 pages
% Feature geometry + identifiability              1.50 pages
% Gauge construction + protected singularity      1.25 pages
% Ratio approximations and finite-fit corrections 0.50 pages
% Experiments                                     1.50 pages
% Discussion + limitations                        0.75 pages
% Target total                                    9.00 pages
% -----------------------------------------------------------------------------

\newtheorem{theorem}{Theorem}
\newtheorem{proposition}[theorem]{Proposition}
\newtheorem{corollary}[theorem]{Corollary}

\providecommand{\E}{\mathbb{E}}
\providecommand{\R}{\mathbb{R}}
\providecommand{\tr}{\operatorname{tr}}
\newcommand{\bhat}{\widehat{\bm\beta}}
\newcommand{\bstar}{\bm\beta^{\star}}
\newcommand{\tstar}{\bm w^{\star}}
\newcommand{\Egen}{E_{\rm gen}}
\newcommand{\Eacc}{E_{\rm acc}}
\newcommand{\Rgen}{R_{\rm gen}^{2}}
\newcommand{\Racc}{R_{\rm acc}^{2}}
\newcommand{\Sperp}{S_{\perp}}

\title{Prediction Does Not Imply Control:\\
Geometry and Random-Matrix Theory of Self-Generated Distribution Shift}
\author{Anonymous Authors}

\begin{document}
\maketitle

\begin{abstract}
\begin{itemize}[leftmargin=1.25em,itemsep=0pt,topsep=0pt]
  \item Predictive neural models with nearly identical held-out accuracy can
        differ dramatically when used to synthesize stimuli that control the
        true response.
  \item We isolate this failure in a linear teacher--student model: natural
        prediction is evaluated in the covariance geometry, whereas
        accentuation probes a direction selected by the fitted model itself.
  \item Deterministic equivalents quantify prediction, accentuation, and
        cross-validated ridge under anisotropic data.
  \item With fixed linear features, prediction depends on fewer summary
        statistics than pixel-space control. We construct continuous families
        with identical prediction and CV behavior but arbitrarily different
        control.
  \item Natural-image covariance experiments validate the theory and recover
        the high-frequency failures observed in neuronal accentuation.
\end{itemize}
\end{abstract}

\section{Introduction}

\begin{itemize}[leftmargin=1.3em]
  \item \textbf{Empirical puzzle.} Encoding models can predict held-out neural
        responses equally well but generate very different preferred stimuli;
        weak generators may nevertheless recognize stimuli made by a strong
        generator.
  \item \textbf{Conceptual gap.} Prediction uses an externally specified
        natural distribution. Control induces a self-generated distribution
        whose direction is chosen from the fitted model's own gradient.
  \item \textbf{Minimal explanation.} Even linear regression exhibits the
        full dissociation. Covariance-weighted coefficient error governs
        prediction, while unweighted pixel-gradient geometry governs control.
  \item \textbf{Technical contribution.} We derive high-dimensional
        deterministic equivalents (DEs) for natural generalization,
        accentuation, peer evaluation, and population-optimal CV ridge.
  \item \textbf{Strongest result.} In linear feature regression, the complete
        Gaussian training distribution and prediction-optimal ridge path do
        not identify pixel-space control. A prediction-preserving rotation can
        continuously move control $R^2$ from near one to highly negative.
  \item \textbf{Paper map.} First establish exact metrics, then pixel-ridge
        RMT, then feature-space non-identifiability, and finally empirical
        validation on natural-image spectra.
\end{itemize}

\begin{figure}[t]
  \centering
  \fbox{\parbox{0.94\linewidth}{
    \textbf{Figure 1: motivating phenomenon and geometry.}
    (a) Similar natural-image prediction, different generated stimuli and
    control. (b) Peer-review asymmetry. (c) Natural inputs probe a covariance
    ellipsoid; accentuation moves along the learned pixel gradient.
  }}
  \caption{Prediction and control evaluate the same fitted function on
  different, model-dependent distributions.}
  \label{fig:motivation}
\end{figure}

\paragraph{Related work (compressed).}
\begin{itemize}[leftmargin=1.3em,itemsep=1pt,topsep=2pt]
  \item Model-based neuronal control and feature visualization. [CITATIONS]
  \item Benign overfitting, anisotropic ridge, and deterministic equivalents.
        [CITATIONS]
  \item Distribution shift, performative prediction, and optimization against
        learned predictors. [CITATIONS]
\end{itemize}

\section{Setup: three evaluations of one predictor}
\label{sec:setup}

\begin{itemize}[leftmargin=1.3em]
  \item Let
  \begin{equation}
    \bm x\sim\mathcal N(0,\Sigma),\qquad
    y=(\bstar)^{\top}\bm x+\sigma\epsilon,\qquad
    S=(\bstar)^{\top}\Sigma\bstar,
    \label{eq:teacher}
  \end{equation}
  where $S$ is the noiseless signal variance.
  \item Separate the \emph{evaluation path} from the fitted model. For an
        evaluator $\bhat_e$ and generator direction $\bhat_g$, define
  \begin{equation}
    T_g=(\bstar)^{\top}\bhat_g,\qquad
    D_{e\leftarrow g}=\bhat_e^{\top}\bhat_g.
    \label{eq:peer-gains}
  \end{equation}
  Along $\bm x(t)=\bm x_0+t\bhat_g$, the teacher and evaluator gains are
  $T_g$ and $D_{e\leftarrow g}$.
\end{itemize}

\begin{table}[t]
\caption{Evaluation metrics. Pathwise MSE is reported per unit $\operatorname{Var}(t)$;
the own-path normalization $\operatorname{Var}(t)=S/D^2$ yields
$\Eacc/S=(1-N/D)^2$.}
\label{tab:metrics}
\centering
\scriptsize
\resizebox{\textwidth}{!}{%
\begin{tabular}{@{}lllll@{}}
\toprule
Evaluation & Distribution / path & Error & $R^2$ & true-on-predicted slope \\
\midrule
Natural generalization
& $\bm x\sim\mathcal N(0,\Sigma)$
& $(\bhat-\bstar)^\top\Sigma(\bhat-\bstar)$
& $1-\Egen/S$
& $\bhat^\top\Sigma\bstar/(\bhat^\top\Sigma\bhat)$ \\
Own accentuation
& $\bm x_0+t\bhat$; $N=(\bstar)^\top\bhat$, $D=\bhat^\top\bhat$
& $\operatorname{Var}(t)(D-N)^2$
& $1-(1-D/N)^2$
& $N/D$ \\
Peer evaluation
& $\bm x_0+t\bhat_g$ evaluated by $\bhat_e$
& $\operatorname{Var}(t)(D_{e\leftarrow g}-T_g)^2$
& $1-(1-D_{e\leftarrow g}/T_g)^2$
& $T_g/D_{e\leftarrow g}$ \\
\bottomrule
\end{tabular}}
\end{table}

\begin{itemize}[leftmargin=1.3em]
  \item The table exposes two distinct singularities. The normalized
        accentuation MSE diverges when $D\to0$; accentuation $R^2$ becomes very
        negative when the true gain $N\to0$ while $D$ stays nonzero.
  \item Very negative $\Racc$ therefore need not imply an unbounded raw error:
        as $N\to0$, $\Racc\to-\infty$ but the student-range-normalized
        $\Eacc/S\to1$.
\end{itemize}

\section{Pixel ridge: one short deterministic equivalent}
\label{sec:pixel}

\begin{itemize}[leftmargin=1.3em]
  \item Fit normalized ridge from $n$ samples,
  \begin{equation}
    \bhat=(\widehat\Sigma+\lambda I)^{-1}
      \left(\widehat\Sigma\bstar+\frac{\sigma}{n}X^\top\bm\eta\right),
    \qquad \widehat\Sigma=X^\top X/n.
    \label{eq:ridge}
  \end{equation}
  Software using $\lVert\bm y-X\bm\beta\rVert^2+\alpha\lVert\bm\beta\rVert^2$
  has $\alpha=n\lambda$.
  \item Let $\kappa$ solve
  \begin{equation}
    \kappa-\lambda=\frac{\kappa}{n}
      \tr\!\left[\Sigma(\Sigma+\kappa I)^{-1}\right],
    \label{eq:kappa}
  \end{equation}
  and abbreviate
  \begin{equation}
    B_{p,q}=(\bstar)^\top\Sigma^p(\Sigma+\kappa I)^{-q}\bstar,
    \qquad d_{p,q}=\tr[\Sigma^p(\Sigma+\kappa I)^{-q}],
    \quad d_2=d_{2,2}.
    \label{eq:summaries}
  \end{equation}
\end{itemize}

\begin{proposition}[Pixel-space DE and CV consequence]
\label{prop:pixel}
In the proportional limit,
\begin{align}
  \Egen^{\rm DE}
  &=\frac{n}{n-d_2}(\kappa^2B_{1,2}+\sigma^2)-\sigma^2,
  \label{eq:egen-de}\\
  z_{\rm acc}^{\rm DE}
  &:=\frac{D_{\rm acc}}{N_{\rm acc}}-1
  \asymp
  \frac{\sigma^2d_{1,2}-n\lambda B_{1,2}}
       {(n-d_2)B_{1,1}}.
  \label{eq:z-de}
\end{align}
At an interior minimizer of $\Egen^{\rm DE}$,
\begin{align}
  (n-d_2)\kappa B_{2,3}
  &=d_{2,3}(\kappa^2B_{1,2}+\sigma^2),
  \label{eq:cv-condition}\\
  z_{\rm acc,CV}^{\rm DE}
  &=\kappa\frac{d_{1,2}}{B_{1,1}}
    \left(\frac{B_{2,3}}{d_{2,3}}-
                \frac{B_{1,2}}{d_{1,2}}\right),
  \label{eq:cv-z}
\end{align}
which is not generically zero.
\end{proposition}

\begin{itemize}[leftmargin=1.3em]
  \item \textbf{Interpretation.} CV optimizes covariance-weighted prediction;
        it has no first-order condition forcing the self-generated path to be
        calibrated.
  \item \textbf{Spectral mechanism.} Prediction downweights coefficient error
        in low-variance directions, while $D_{\rm acc}=\lVert\bhat\rVert^2$
        counts that energy without covariance weighting.
  \item \textbf{What moves to appendix.} Resolvent identities, one-/two-point
        DE derivations, boundary CV optima, and randomness of fold-selected
        $\alpha$.
\end{itemize}

\begin{figure}[t]
  \centering
  \fbox{\parbox{0.94\linewidth}{
    \textbf{Figure 2: pixel-ridge theory and validation.}
    (a) Natural-image eigenspectrum and disk-teacher power. (b) Per-PC error.
    (c) Dense noise sweep: $\Rgen$ versus $\Racc$, slope, and CV $\alpha$ on
    vertically aligned axes. (d) DE versus Monte Carlo; optional peer panel.
  }}
  \caption{Cross-validation protects natural prediction but does not enforce
  control calibration.}
  \label{fig:pixel}
\end{figure}

\section{Linear features separate prediction from pixel control}
\label{sec:features}

\begin{itemize}[leftmargin=1.3em]
  \item Fit $\widehat f(\bm x)=\widehat{\bm w}^{\top}F\bm x$ in feature
        coordinates; accentuation back-propagates the pixel weight
        $\bhat=F^\top\widehat{\bm w}$.
  \item Define two pairs of geometries:
  \begin{equation}
    \underbrace{C=F\Sigma F^\top,\quad \bm h=F\Sigma\bstar}_{\text{natural prediction}},
    \qquad
    \underbrace{G=FF^\top,\quad \bm q=F\bstar}_{\text{pixel control}}.
    \label{eq:two-geometries}
  \end{equation}
  \item The realizable feature teacher and omitted signal power are
  \begin{equation}
    \tstar=C^\dagger\bm h,
    \qquad \Sperp=S-\bm h^\top C^\dagger\bm h.
    \label{eq:feature-teacher}
  \end{equation}
\end{itemize}

\begin{table}[t]
\caption{Pixel and feature regression have parallel prediction DEs, but
feature-space accentuation requires extra pixel-geometry statistics.}
\label{tab:parallel}
\centering
\small
\begin{tabular}{@{}lll@{}}
\toprule
Role & Pixel regression & Feature regression \\
\midrule
Covariance & $\Sigma$ & $C=F\Sigma F^\top$ \\
Realizable teacher & $\bstar$ & $\tstar=C^\dagger F\Sigma\bstar$ \\
Unrealizable signal & $0$ & $\Sperp$ \\
Effective training noise & $\sigma^2$ & $\sigma^2+\Sperp$ \\
Back-projection geometry & $I$ & $G=FF^\top$ \\
Teacher path gain & $\bstar$ & $\bm q=F\bstar$ \\
\bottomrule
\end{tabular}
\end{table}

\begin{itemize}[leftmargin=1.3em]
  \item Define $B^F_{p,q}$ and $d^F_{p,q}$ by replacing
        $(\Sigma,\bstar)$ in Eq.~\eqref{eq:summaries} with $(C,\tstar)$.
        The prediction formula is literally parallel:
  \begin{equation}
    \Egen^{F,{\rm DE}}
    =\frac{n}{n-d_2^F}
      \left(\kappa^2B^F_{1,2}+\Sperp+\sigma^2\right)-\sigma^2.
    \label{eq:feature-egen}
  \end{equation}
  \item Define the mean shrunken pixel weight
  \begin{equation}
    \widetilde{\bm\beta}_\kappa
    =F^\top(C+\kappa I)^{-1}F\Sigma\bstar.
    \label{eq:betatilde}
  \end{equation}
  The accentuation moments expose the extra freedom:
  \begin{align}
    N_{\rm acc}^{F,{\rm DE}}
    &\asymp (\bstar)^\top\widetilde{\bm\beta}_\kappa,
    \label{eq:feature-n}\\
    D_{\rm acc}^{F,{\rm DE}}
    &\asymp \lVert\widetilde{\bm\beta}_\kappa\rVert^2
    +\frac{\Egen^{F,{\rm DE}}+\sigma^2}{n}
      \tr\!\left[G(C+\kappa I)^{-2}C\right].
    \label{eq:feature-d}
  \end{align}
\end{itemize}

\begin{theorem}[Prediction does not identify pixel control]
\label{thm:separation}
For Gaussian inputs and deterministic $F$, the training law, natural
generalization curve, and population/DE-CV choice of ridge are determined by
$(C,\bm h,S,\sigma^2,n)$. Pixel-gradient accentuation additionally depends on
$(G,\bm q)=(FF^\top,F\bstar)$. Thus identical prediction and CV behavior does
not imply identical $\Eacc$, accentuation slope, or $\Racc$.
\end{theorem}

\begin{itemize}[leftmargin=1.3em]
  \item This is the central thesis in statistical form: even conditioning on
        the complete prediction problem, $G$ and $\bm q$ retain degrees of
        freedom that can change control.
  \item For non-Gaussian natural images, make the narrower statement that
        $(C,\bm h,S)$ identify the second-order prediction theory, not
        necessarily the full feature--response distribution.
\end{itemize}

\section{Prediction-preserving changes can destroy control}
\label{sec:gauge}

\begin{itemize}[leftmargin=1.3em]
  \item Whiten the input geometry:
  \begin{equation}
    A=F\Sigma^{1/2},\qquad
    \bm b=\Sigma^{1/2}\bstar,\qquad
    \bm c=\Sigma^{-1/2}\bstar.
    \label{eq:whitened}
  \end{equation}
  Then $C=AA^\top$, $\bm h=A\bm b$, $\bm q=A\bm c$, and
  $G=A\Sigma^{-1}A^\top$.
\end{itemize}

\begin{theorem}[Prediction-preserving gauge]
\label{thm:gauge}
Let $Q\in O(d)$ satisfy $Q\bm b=\bm b$ and set
\begin{equation}
  F_Q=AQ\Sigma^{-1/2}.
  \label{eq:gauge}
\end{equation}
Then $F_Q\Sigma F_Q^\top=C$ and $F_Q\Sigma\bstar=\bm h$; hence every member
has the same Gaussian training law, $\Egen(\lambda)$, and CV ridge choice.
Nevertheless,
\begin{equation}
  G_Q=AQ\Sigma^{-1}Q^\top A^\top,
  \qquad \bm q_Q=AQ\bm c
  \label{eq:gauge-control}
\end{equation}
can vary, so accentuation can vary. Requiring also $Q\bm c=\bm c$ fixes
$\bm q$ and isolates the effect of $G$.
\end{theorem}

\begin{itemize}[leftmargin=1.3em]
  \item Give an intuitive continuous implementation:
  \begin{equation}
    A(t)=P+\sqrt{1-t}\,R+\sqrt t\,V,\qquad
    F(t)=A(t)\Sigma^{-1/2},\qquad t=\sin^2\theta.
    \label{eq:continuous-family}
  \end{equation}
  Here $P$ carries the components needed to fix $A\bm b$ (and optionally
  $A\bm c$); $R$ and $V$ are equal-Gram row frames in the orthogonal null
  space, chosen respectively from well-conditioned and low-variance pixel
  directions. Orthogonality removes cross terms, so $AA^\top$ stays fixed.
  \item One endpoint is a well-behaved top-PC projection; the other loads the
        same prediction-null feature energy into the spectral tail. The
        training problem is unchanged, but $F(t)F(t)^\top$ grows sharply.
  \item State the protected-prediction singularity in one line. At fixed
        $(C,\bm h)$, forcing the mean true path gain to zero has minimum
        prediction cost
  \begin{equation}
    \epsilon_{N=0}
      =\frac{(\bm q^\top\tstar)^2}{\bm q^\top C^{-1}\bm q}.
    \label{eq:protected-singularity}
  \end{equation}
  If $\bm q^\top\tstar$ is small or $C^{-1}\bm q$ is large, natural
  prediction is protected while $\Racc$ approaches its singular regime.
\end{itemize}

\begin{figure}[t]
  \centering
  \fbox{\parbox{0.94\linewidth}{
    \textbf{Figure 3: exact non-identifiability construction.}
    Left-to-right dependency graph separating $(C,\bm h)$ from $(G,\bm q)$.
    Continuous $p=500$ rotation from top-PC features to tail-loaded features.
    Show invariant $\Egen$, $\Rgen$, and CV $\alpha$ above; changing
    $\tr(FF^\top)$, slope, $\Eacc$, and $\Racc$ below. Colors denote noise.
    Include representative feature maps at several rotation angles.
  }}
  \caption{A continuous feature transformation leaves the entire prediction
  problem fixed while changing pixel-space control by orders of magnitude.}
  \label{fig:gauge}
\end{figure}

\section{What does a DE predict for a ratio?}
\label{sec:ratio}

\begin{itemize}[leftmargin=1.3em]
  \item Accentuation metrics are nonlinear functions of the random pair
        $(N,D)=(N_{\rm acc},D_{\rm acc})$. In general,
        $\E[g(N,D)]\neq g(\E N,\E D)$. We therefore distinguish three
        approximation levels rather than calling all of them ``the DE.''
  \item \textbf{Level 1: ratio of deterministic equivalents.} The main-text
        results use
  \begin{equation}
    z_{\rm typ}=\frac{\mu_D}{\mu_N}-1,
    \qquad R_{\rm acc,typ}^2=1-z_{\rm typ}^2.
    \label{eq:ratio-of-de}
  \end{equation}
  This predicts a typical fit and is accurate when $(N,D)$ concentrate and
  $|\mu_N|\gg\sqrt{\operatorname{Var}(N)}$.
  \item \textbf{Level 2: perturbative ratio correction.} A second-order
        expansion gives
  \begin{equation}
    \E\!\left[\frac{D}{N}\right]
    \approx \frac{\mu_D}{\mu_N}
      -\frac{\operatorname{Cov}(D,N)}{\mu_N^2}
      +\frac{\mu_D\operatorname{Var}(N)}{\mu_N^3}.
    \label{eq:ratio-delta}
  \end{equation}
  More generally,
  $\E[g(N,D)]\approx g(\bm\mu)+
  \tfrac12\tr[H_g(\bm\mu)\operatorname{Cov}(N,D)]$.
  This corrects modest finite-fit bias and supplies uncertainty estimates,
  but it fails when appreciable probability mass approaches $N=0$.
  \item \textbf{Level 3: Gaussian distributional DE.} Approximate the whole
        fitted feature coefficient by
  \begin{equation}
    \widehat{\bm w}\ \dot\sim\
      \mathcal N\!\left(
      C(C+\kappa I)^{-1}\tstar,
      \frac{\Egen+\sigma^2}{n}(C+\kappa I)^{-1}C(C+\kappa I)^{-1}
      \right).
    \label{eq:gaussian-de}
  \end{equation}
  Pushing this Gaussian through the linear statistic $N$ and quadratic
  statistic $D$ predicts their joint distribution, including medians,
  quantiles, sign changes, and the probability of near-singular fits. It is
  most useful at high noise, smaller feature dimension, and near a control
  transition. Random CV selection can be incorporated as a mixture over the
  selected $\kappa$ values.
  \item \textbf{Denominator crossing.} If the distribution of $N$ has mass
        arbitrarily close to zero, literal moments such as $\E[D/N]$ or
        $\E[(D/N)^2]$ can be unstable or fail to exist. In that regime we
        compare median/IQR and tail probabilities, not only Monte Carlo
        means.
  \item \textbf{Scope of the claims.} The formal results in
        Sections~\ref{sec:pixel}--\ref{sec:gauge} remain at the leading DE,
        typical-fit level. The two corrections refine finite-realization
        predictions; they do not change the exact summary-statistic
        separation or gauge non-identifiability theorem.
\end{itemize}


\section{Experiments}
\label{sec:experiments}
\begin{itemize}[leftmargin=1.3em]
  \item \textbf{Synthetic spectra.} Isotropic and power-law covariance;
        teachers aligned with top, middle, bottom, or random eigendirections.
        Compare DE and Monte Carlo across signal-to-noise ratios and dense CV
        grids.
  \item \textbf{Natural-image covariance.} Van Hateren / FFHQ patches with a
        circular disk teacher. Plot teacher and fitted weights in pixel and PC
        coordinates, including low-variance high-frequency error.
  \item \textbf{Feature baselines.} Pixel regression, PCA coordinates,
        whitened PCA, top-PC truncation, and prediction-preserving rotations.
        Show $\Egen$--$\Eacc$ and $\Rgen$--$\Racc$ gaps on the same axes.
  \item \textbf{Reporting convention.} Every noise plot has lower axis
        response noise $\sigma$ and upper axis noise/signal variance
        $\sigma^2/S$. Cache plot-ready tables separately from rendered figures.
\end{itemize}

\begin{figure}[t]
  \centering
  \fbox{\parbox{0.94\linewidth}{
    \textbf{Figure 4: validation and robustness.}
    (a) DE/MC agreement across spectra and teacher alignments. (b) Natural
    images with disk teacher. (c) PCA, whitening, truncation, and rotated
    feature maps. (d) Random null-space samples showing a broad distribution
    of control at fixed prediction. Insets only where they do not obscure the
    principal curves.
  }}
  \caption{Prediction-equivalent models can span a broad range of
  accentuation behavior under realistic anisotropic spectra.}
  \label{fig:experiments}
\end{figure}

\section{Discussion}

\begin{itemize}[leftmargin=1.3em]
  \item \textbf{Main message.} Good in-distribution prediction constrains the
        model in directions emphasized by natural covariance; control depends
        on the model's behavior along directions that it chooses itself.
  \item \textbf{Why features matter.} A feature map can preserve all
        prediction-relevant second-order statistics while radically changing
        the metric induced by back-propagation to pixels.
  \item \textbf{Practical implication.} Selecting models or ridge penalties
        by held-out prediction alone cannot certify gradient-based control.
        Validation must directly constrain gradients or test generated paths.
  \item \textbf{Biological interpretation.} High-frequency accentuations are
        not merely a visual artifact: they are the natural signature of
        low-variance directions weakly constrained by observed stimuli.
  \item \textbf{Limitations.} Linear teachers, fixed feature maps, Gaussian or
        second-order input theory, and one-dimensional gradient paths. The
        gauge theorem is exact in its stated Gaussian setting; extensions to
        nonlinear learned features remain open.
\end{itemize}

\paragraph{Reproducibility.}
\begin{itemize}[leftmargin=1.3em,itemsep=1pt,topsep=2pt]
  \item Release scripts, spectra, seeds, cached numerical summaries, and exact
        ridge conventions; report the full CV grid and boundary selections.
  \item Verify each long sweep with timing tests, progress logs, and cached
        plot tables so visual changes do not require recomputation.
\end{itemize}

\paragraph{Broader impact.}
\begin{itemize}[leftmargin=1.3em,itemsep=1pt,topsep=2pt]
  \item The work cautions against deploying predictive models as controllers
        without intervention-level validation, especially in neuroscience.
\end{itemize}

% References are intentionally omitted from this bullet skeleton.
% Appendix plan (not rendered here):
% A. Exact metric derivations and normalization conventions.
% B. Pixel one-/two-resolvent deterministic equivalents.
% C. CV stationary calculation and random fold-selection analysis.
% D. Feature-space DE derivation, including noncommuting F F^T insertions.
% E. Gauge theorem, continuous-frame construction, and singular covariance.
% F. Gaussian distributional DE and ratio-moment propagation.
% G. Peer-review formulas and additional experiments.

\end{document}
